\documentclass[12pt,a4paper]{article}
\usepackage[utf8]{inputenc}
\usepackage[spanish]{babel}
\usepackage{amsmath,amssymb}
\usepackage{geometry}
\usepackage{fancyhdr}
\usepackage{titlesec}
\usepackage{enumitem}
\usepackage{parskip}
\usepackage{booktabs}
\usepackage{array}
\usepackage{pgfplots}
\pgfplotsset{compat=1.18}
\usepackage{tikz}
\usepackage{xcolor}
\usepackage{hyperref}
\usepackage{graphicx}
\usepackage{listings}

\lstset{
  basicstyle=\ttfamily\small,
  breaklines=true,
  columns=fullflexible,
  keywordstyle=\color{blue},
  commentstyle=\color{gray},
  showstringspaces=false,
  frame=single
}

\geometry{top=2.5cm, bottom=2.5cm, left=2.5cm, right=2.5cm}

\pagestyle{fancy}
\fancyhf{}
\renewcommand{\headrulewidth}{0.4pt}
\fancyhead[C]{%
  \begin{center}
    \textbf{ESCUELA POLIT\'ECNICA NACIONAL}\\[-2pt]
    \textbf{M\'ETODOS NUM\'ERICOS}\\[-2pt]
    \textbf{GR1CC}
  \end{center}
}
\setlength{\headheight}{52pt}
\fancyfoot[C]{\thepage}

\setlength{\parindent}{0pt}
\setlength{\parskip}{6pt}

\titleformat{\section}[block]{\bfseries\large}{}{0pt}{}
\titleformat{\subsection}[block]{\bfseries}{}{0pt}{}

\begin{document}

\vspace*{0.5cm}

\textbf{Nombre:} Karina C\'ardenas D\'iaz \hfill \textbf{Fecha:} 13/07/2026

\vspace{0.3cm}
\begin{center}
  \textbf{\large TAREA 08 --- AJUSTE POR M\'INIMOS CUADRADOS}\\[2pt]
  \textbf{Unidad 03-C}
\end{center}
\vspace{0.2cm}

\hrule
\vspace{0.4cm}

\section*{CONJUNTO DE EJERCICIOS}

\textbf{1.} Dados los datos:

\begin{center}
\begin{tabular}{c|cccccccccc}
\toprule
$x_i$ & 4.0 & 4.2 & 4.5 & 4.7 & 5.1 & 5.5 & 5.9 & 6.3 & 6.8 & 7.1 \\
$y_i$ & 102.56 & 130.11 & 113.18 & 142.05 & 167.53 & 195.14 & 224.87 & 256.73 & 299.50 & 326.72 \\
\bottomrule
\end{tabular}
\end{center}

\textbf{Descripci\'on:} Para cada modelo se plantea $E=\sum_{i=1}^{n}[y_i-f(x_i)]^2$ y se resuelve el sistema de ecuaciones normales derivado de minimizar $E$. Para los modelos no polinomiales ($be^{ax}$ y $bx^a$) se linealiza tomando logaritmo natural, se ajusta una recta sobre los datos transformados y se calcula $E$ sobre la escala original.

\textbf{a. Grado 1:} Sistema normal $2\times 2$ con $S_x,\,S_y,\,S_{xx},\,S_{xy}$:
\[
\boxed{P_1(x) = -191.5724 + 71.6102\,x}, \qquad E = 1058.84
\]

\textbf{b. Grado 2:} Sistema normal $3\times 3$:
\[
\boxed{P_2(x) = 51.0008 - 19.3086\,x + 8.2171\,x^2}, \qquad E = 551.66
\]

\textbf{c. Grado 3:} Sistema normal $4\times 4$:
\[
\boxed{P_3(x) = 469.1633 - 254.8748\,x + 51.5610\,x^2 - 2.6068\,x^3}, \qquad E = 518.38
\]

\textbf{d. Forma $be^{ax}$:} Con $\ln y_i = \ln b + a x_i$ se ajusta la recta sobre $(x_i,\ln y_i)$:
\[
\boxed{f(x) = 24.7767\,e^{0.3685x}}, \qquad E = 821.01
\]

\textbf{e. Forma $bx^a$:} Con $\ln y_i = \ln b + a \ln x_i$ se ajusta la recta sobre $(\ln x_i,\ln y_i)$:
\[
\boxed{f(x) = 6.5187\,x^{1.9933}}, \qquad E = 581.56
\]

\begin{center}
\includegraphics[width=0.85\textwidth]{ejercicio1.png}
\end{center}

El menor error corresponde al polinomio de grado 3, aunque el ajuste de grado 2 y la forma $bx^a$ tambi\'en describen razonablemente bien la tendencia creciente de los datos.

\vspace{0.5cm}
\textbf{2.} Repita el ejercicio 1 para los siguientes datos:

\begin{center}
\begin{tabular}{c|cccccccc}
\toprule
$x_i$ & 0.2 & 0.3 & 0.6 & 0.9 & 1.1 & 1.3 & 1.4 & 1.6 \\
$y_i$ & 0.050446 & 0.098426 & 0.33277 & 0.72660 & 1.0972 & 1.5697 & 1.8487 & 2.5015 \\
\bottomrule
\end{tabular}
\end{center}

\textbf{a. Grado 1:}
\[
\boxed{P_1(x) = -0.5125 + 1.6655\,x}, \qquad E = 0.3356
\]

\textbf{b. Grado 2:}
\[
\boxed{P_2(x) = 0.0851 - 0.3114\,x + 1.1294\,x^2}, \qquad E = 0.002420
\]

\textbf{c. Grado 3:}
\[
\boxed{P_3(x) = -0.0184 + 0.2484\,x + 0.4029\,x^2 + 0.2662\,x^3}, \qquad E = 5.07\times10^{-6}
\]

\textbf{d. Forma $be^{ax}$:}
\[
\boxed{f(x) = 0.0457\,e^{2.7073x}}, \qquad E = 1.0750
\]

\textbf{e. Forma $bx^a$:}
\[
\boxed{f(x) = 0.9502\,x^{1.8720}}, \qquad E = 0.05448
\]

\begin{center}
\includegraphics[width=0.85\textwidth]{ejercicio2.png}
\end{center}

En este caso el polinomio de grado 3 pr\'acticamente interpola los datos (error casi nulo), lo que sugiere que la relaci\'on subyacente es muy cercana a un polinomio c\'ubico.

\vspace{0.5cm}
\textbf{3.} La tabla muestra los promedios de puntos del colegio de 20 estudiantes junto con su puntuaci\'on ACT de matem\'aticas. Grafique los datos y encuentre la recta por m\'inimos cuadrados.

\textbf{Descripci\'on:} Se aplica el sistema normal de la recta $y=a_1x+a_0$ usando $S_x=\sum x_i,\ S_y=\sum y_i,\ S_{xx}=\sum x_i^2,\ S_{xy}=\sum x_iy_i$, con $n=20$.

\textbf{Respuesta:}
\[
\boxed{y = 0.1009\,x + 0.4866}, \qquad E = 5.0487
\]

\begin{center}
\includegraphics[width=0.85\textwidth]{ejercicio3.png}
\end{center}

La pendiente positiva indica una relaci\'on directa (aunque d\'ebil) entre la puntuaci\'on ACT y el promedio de calificaciones; la dispersi\'on de los puntos respecto a la recta explica el valor relativamente alto del error.

\vspace{0.5cm}
\textbf{4.} El conjunto de datos presentado al Subcomit\'e Antimonopolio del Senado muestra el peso promedio de distintos tipos de veh\'iculo y el porcentaje de accidentes en los que la lesi\'on m\'as grave fue fatal o seria. Encuentre la recta por m\'inimos cuadrados.

\textbf{Respuesta:}
\[
\boxed{y = -0.0022550\,x + 13.1465}, \qquad E = 2.0591
\]

\begin{center}
\includegraphics[width=0.85\textwidth]{ejercicio4.png}
\end{center}

La pendiente negativa confirma la relaci\'on esperada: a mayor peso promedio del veh\'iculo, menor el porcentaje de lesiones graves o fatales en caso de accidente.

\vspace{0.3cm}
\textbf{Repositorio:} \url{https://github.com/usuario/tarea08-minimos-cuadrados}

\end{document}
